\documentclass[12pt,a4paper]{article}
\usepackage[UTF8]{ctex}
\usepackage{amsmath,amssymb,amsthm}
\usepackage{geometry}
\usepackage{bm}

\geometry{margin=2.5cm}

\title{从向量变换推导时间导数关系的详细证明}
\author{第8章补充说明}
\date{}

\begin{document}

\maketitle

\section{问题提出}

\textbf{已知}：
\begin{enumerate}
\item 向量变换关系：
\begin{equation}
\bm{v}_I = R(t) \bm{v}_b
\label{eq:vector_transform}
\end{equation}

其中$R(t)$是从体坐标系$\{b\}$到惯性坐标系$\{I\}$的旋转矩阵。

\item 旋转矩阵的导数关系：
\begin{equation}
\dot{R}(t) = R(t) [\omega_b]_{\times}
\label{eq:rotation_derivative}
\end{equation}

其中$\omega_b$是体坐标系相对于惯性坐标系的角速度（在体坐标系中表示），$[\omega_b]_{\times}$是$\omega_b$的反对称矩阵。

\textbf{要证明}：
\begin{equation}
\left.\frac{d\bm{v}}{dt}\right|_{\text{inertial}} = \left.\frac{d\bm{v}}{dt}\right|_{\text{body}} + \omega_b \times \bm{v}_b
\label{eq:target}
\end{equation}
\end{enumerate}

\section{推导方法1：直接对向量变换求导}

\subsection{步骤1：对向量变换关系求导}

对式(\ref{eq:vector_transform})两边对时间求导：
\begin{align}
\frac{d\bm{v}_I}{dt} &= \frac{d}{dt}[R(t) \bm{v}_b] \\
&= \dot{R}(t) \bm{v}_b + R(t) \dot{\bm{v}}_b
\label{eq:derivative_expanded}
\end{align}

\textbf{说明}：
\begin{itemize}
\item $\frac{d\bm{v}_I}{dt}$：向量在惯性坐标系中的时间导数
\item $\dot{R}(t)$：旋转矩阵的时间导数
\item $\dot{\bm{v}}_b$：向量在体坐标系中的时间导数
\end{itemize}

\subsection{步骤2：使用旋转矩阵的导数关系}

将式(\ref{eq:rotation_derivative})代入式(\ref{eq:derivative_expanded})：
\begin{align}
\frac{d\bm{v}_I}{dt} &= R(t) [\omega_b]_{\times} \bm{v}_b + R(t) \dot{\bm{v}}_b \\
&= R(t) ([\omega_b]_{\times} \bm{v}_b + \dot{\bm{v}}_b)
\label{eq:intermediate}
\end{align}

\subsection{步骤3：使用反对称矩阵的性质}

反对称矩阵$[\omega_b]_{\times}$的性质：
\begin{equation}
[\omega_b]_{\times} \bm{v}_b = \omega_b \times \bm{v}_b
\label{eq:skew_property}
\end{equation}

\textbf{证明}：

设$\omega_b = \begin{bmatrix} \omega_x \\ \omega_y \\ \omega_z \end{bmatrix}$，$\bm{v}_b = \begin{bmatrix} v_x \\ v_y \\ v_z \end{bmatrix}$，则：
\begin{equation}
[\omega_b]_{\times} = \begin{bmatrix}
0 & -\omega_z & \omega_y \\
\omega_z & 0 & -\omega_x \\
-\omega_y & \omega_x & 0
\end{bmatrix}
\end{equation}

因此：
\begin{align}
[\omega_b]_{\times} \bm{v}_b &= \begin{bmatrix}
0 & -\omega_z & \omega_y \\
\omega_z & 0 & -\omega_x \\
-\omega_y & \omega_x & 0
\end{bmatrix} \begin{bmatrix}
v_x \\ v_y \\ v_z
\end{bmatrix} \\
&= \begin{bmatrix}
-\omega_z v_y + \omega_y v_z \\
\omega_z v_x - \omega_x v_z \\
-\omega_y v_x + \omega_x v_y
\end{bmatrix} \\
&= \omega_b \times \bm{v}_b
\end{align}

\textbf{证毕}。

\subsection{步骤4：代入反对称矩阵性质}

将式(\ref{eq:skew_property})代入式(\ref{eq:intermediate})：
\begin{align}
\frac{d\bm{v}_I}{dt} &= R(t) (\omega_b \times \bm{v}_b + \dot{\bm{v}}_b) \\
&= R(t) \dot{\bm{v}}_b + R(t) (\omega_b \times \bm{v}_b)
\label{eq:before_transform}
\end{align}

\subsection{步骤5：转换到体坐标系表示}

\textbf{关键观察}：式(\ref{eq:before_transform})的结果是在惯性坐标系$\{I\}$中的表示。要得到在体坐标系$\{B\}$中的表示，需要左乘$R^T(t)$。

\textbf{左乘$R^T(t)$}：
\begin{align}
R^T(t) \frac{d\bm{v}_I}{dt} &= R^T(t) R(t) \dot{\bm{v}}_b + R^T(t) R(t) (\omega_b \times \bm{v}_b) \\
&= I \dot{\bm{v}}_b + I (\omega_b \times \bm{v}_b) \\
&= \dot{\bm{v}}_b + \omega_b \times \bm{v}_b
\label{eq:body_frame_result}
\end{align}

其中$I$是$3 \times 3$单位矩阵。

\subsection{步骤6：解释结果}

\textbf{物理意义}：

\begin{itemize}
\item $R^T(t) \frac{d\bm{v}_I}{dt}$：向量在惯性系中的时间导数，转换到体坐标系中的表示

\item $\dot{\bm{v}}_b$：向量在体坐标系中观察到的变化率，即$\left.\frac{d\bm{v}}{dt}\right|_{\text{body}}$

\item $\omega_b \times \bm{v}_b$：由于坐标系旋转产生的额外项
\end{itemize}

\textbf{因此}，在体坐标系中，向量在惯性系中的时间导数表示为：
\begin{equation}
\left.\frac{d\bm{v}}{dt}\right|_{\text{inertial}} = \left.\frac{d\bm{v}}{dt}\right|_{\text{body}} + \omega_b \times \bm{v}_b
\end{equation}

\textbf{证毕}。

\section{推导方法2：使用旋转矩阵的指数表示}

\subsection{旋转矩阵的指数表示}

旋转矩阵可以表示为：
\begin{equation}
R(t) = \exp([\theta(t)]_{\times})
\end{equation}

其中$\theta(t)$是旋转角度向量，满足$\dot{\theta}(t) = \omega(t)$。

\subsection{对时间求导}

\begin{align}
\dot{R}(t) &= \frac{d}{dt}[\exp([\theta(t)]_{\times})] \\
&= \exp([\theta(t)]_{\times}) [\dot{\theta}(t)]_{\times} \\
&= R(t) [\omega(t)]_{\times}
\end{align}

这与前面的结果一致。

\subsection{应用到向量导数}

\begin{align}
\frac{d}{dt}[R(t) \bm{v}_b] &= \dot{R}(t) \bm{v}_b + R(t) \dot{\bm{v}}_b \\
&= R(t) [\omega_b]_{\times} \bm{v}_b + R(t) \dot{\bm{v}}_b \\
&= R(t) (\omega_b \times \bm{v}_b + \dot{\bm{v}}_b)
\end{align}

在体坐标系中表示：
\begin{equation}
R^T(t) \frac{d}{dt}[R(t) \bm{v}_b] = \omega_b \times \bm{v}_b + \dot{\bm{v}}_b
\end{equation}

\section{推导方法3：几何直观方法}

\subsection{固定向量的情况}

考虑一个固定在体坐标系中的向量$\bm{v}_b$（在体坐标系中不变）。

\textbf{在体坐标系中}：
\begin{equation}
\dot{\bm{v}}_b = 0
\end{equation}

\textbf{在惯性坐标系中}：

由于坐标系旋转，向量在惯性系中的方向在改变。在时间$\Delta t$内，坐标系旋转了$\omega_b \Delta t$的角度。

\textbf{向量变化}：
\begin{equation}
\Delta \bm{v}_I \approx \omega_b \times \bm{v}_b \cdot \Delta t
\end{equation}

\textbf{因此}：
\begin{equation}
\frac{d\bm{v}_I}{dt} = \lim_{\Delta t \to 0} \frac{\Delta \bm{v}_I}{\Delta t} = \omega_b \times \bm{v}_b
\end{equation}

\subsection{变化向量的情况}

如果向量不仅在坐标系中变化，坐标系本身也在旋转，那么总的变化率是两项之和：
\begin{equation}
\frac{d\bm{v}_I}{dt} = \dot{\bm{v}}_b + \omega_b \times \bm{v}_b
\end{equation}

\section{关键步骤详解}

\subsection{为什么需要左乘$R^T(t)$？}

\textbf{问题}：为什么在步骤5中需要左乘$R^T(t)$？

\textbf{回答}：

\begin{enumerate}
\item $\frac{d\bm{v}_I}{dt}$是在惯性坐标系$\{I\}$中的表示

\item 我们需要得到在体坐标系$\{B\}$中的表示

\item 从$\{I\}$到$\{B\}$的坐标变换是$R^T(t)$（因为从$\{B\}$到$\{I\}$是$R(t)$）

\item 因此，要将在$\{I\}$中的向量转换到$\{B\}$中，需要左乘$R^T(t)$
\end{enumerate}

\textbf{数学表达}：

如果$\bm{v}_I = R(t) \bm{v}_b$，那么：
\begin{equation}
\bm{v}_b = R^T(t) \bm{v}_I
\end{equation}

类似地，对于时间导数：
\begin{equation}
\left.\frac{d\bm{v}}{dt}\right|_{\text{body}} = R^T(t) \left.\frac{d\bm{v}}{dt}\right|_{\text{inertial}}
\end{equation}

\subsection{为什么$\omega_b \times \bm{v}_b$项出现在体坐标系中？}

\textbf{关键观察}：

虽然$\omega_b \times \bm{v}_b$是在体坐标系中计算的，但它表示的是由于坐标系旋转产生的变化，这个变化在体坐标系中也能观察到。

\textbf{物理意义}：

\begin{itemize}
\item 即使向量在体坐标系中不变（$\dot{\bm{v}}_b = 0$），由于坐标系旋转，向量在惯性系中的方向在改变
\item 这个方向变化在体坐标系中表现为$\omega_b \times \bm{v}_b$项
\item 这是科里奥利效应的体现
\end{itemize}

\section{验证：特殊情况}

\subsection{情况1：向量在体坐标系中不变}

如果$\dot{\bm{v}}_b = 0$（向量在体坐标系中不变），那么：
\begin{equation}
\left.\frac{d\bm{v}}{dt}\right|_{\text{inertial}} = \omega_b \times \bm{v}_b
\end{equation}

\textbf{验证}：

从向量变换：
\begin{align}
\frac{d\bm{v}_I}{dt} &= \dot{R}(t) \bm{v}_b + R(t) \cdot 0 \\
&= R(t) [\omega_b]_{\times} \bm{v}_b \\
&= R(t) (\omega_b \times \bm{v}_b)
\end{align}

在体坐标系中：
\begin{equation}
R^T(t) \frac{d\bm{v}_I}{dt} = \omega_b \times \bm{v}_b
\end{equation}

\textbf{符合预期}。

\subsection{情况2：坐标系不旋转}

如果$\omega_b = 0$（坐标系不旋转），那么：
\begin{equation}
\left.\frac{d\bm{v}}{dt}\right|_{\text{inertial}} = \left.\frac{d\bm{v}}{dt}\right|_{\text{body}}
\end{equation}

\textbf{验证}：

从向量变换：
\begin{align}
\frac{d\bm{v}_I}{dt} &= R(t) \cdot 0 \cdot \bm{v}_b + R(t) \dot{\bm{v}}_b \\
&= R(t) \dot{\bm{v}}_b
\end{align}

在体坐标系中：
\begin{equation}
R^T(t) \frac{d\bm{v}_I}{dt} = \dot{\bm{v}}_b
\end{equation}

\textbf{符合预期}。

\section{应用到角动量}

\subsection{角动量在体坐标系中的表示}

角动量$\bm{L}_b = I_b \omega_b$在体坐标系$\{b\}$中表示。

\subsection{在惯性系中的时间导数}

根据时间导数关系：
\begin{equation}
\left.\frac{d\bm{L}_b}{dt}\right|_{\text{inertial}} = \left.\frac{d\bm{L}_b}{dt}\right|_{\text{body}} + \omega_b \times \bm{L}_b
\end{equation}

\subsection{展开}

\begin{align}
\left.\frac{d\bm{L}_b}{dt}\right|_{\text{inertial}} &= \frac{d}{dt}(I_b \omega_b) + \omega_b \times (I_b \omega_b) \\
&= I_b \dot{\omega}_b + \frac{dI_b}{dt} \omega_b + \omega_b \times (I_b \omega_b)
\end{align}

由于$I_b$在体坐标系中是常数，$\frac{dI_b}{dt} = 0$。

\textbf{因此}：
\begin{equation}
\left.\frac{d\bm{L}_b}{dt}\right|_{\text{inertial}} = I_b \dot{\omega}_b + \omega_b \times (I_b \omega_b)
\end{equation}

\subsection{角动量定理}

根据角动量定理：
\begin{equation}
m_b = \left.\frac{d\bm{L}_b}{dt}\right|_{\text{inertial}}
\end{equation}

\textbf{因此}：
\begin{equation}
m_b = I_b \dot{\omega}_b + \omega_b \times (I_b \omega_b)
\end{equation}

这就是欧拉方程。

\section{总结}

\subsection{推导步骤总结}

\begin{enumerate}
\item \textbf{向量变换}：$\bm{v}_I = R(t) \bm{v}_b$

\item \textbf{对时间求导}：$\frac{d\bm{v}_I}{dt} = \dot{R}(t) \bm{v}_b + R(t) \dot{\bm{v}}_b$

\item \textbf{使用旋转矩阵导数}：$\dot{R}(t) = R(t) [\omega_b]_{\times}$

\item \textbf{代入并简化}：$\frac{d\bm{v}_I}{dt} = R(t) ([\omega_b]_{\times} \bm{v}_b + \dot{\bm{v}}_b)$

\item \textbf{使用反对称矩阵性质}：$[\omega_b]_{\times} \bm{v}_b = \omega_b \times \bm{v}_b$

\item \textbf{转换到体坐标系}：左乘$R^T(t)$得到体坐标系中的表示

\item \textbf{最终结果}：$\left.\frac{d\bm{v}}{dt}\right|_{\text{inertial}} = \left.\frac{d\bm{v}}{dt}\right|_{\text{body}} + \omega_b \times \bm{v}_b$
\end{enumerate}

\subsection{关键要点}

\begin{itemize}
\item 向量变换关系是基础：$\bm{v}_I = R(t) \bm{v}_b$
\item 旋转矩阵的导数关系：$\dot{R}(t) = R(t) [\omega_b]_{\times}$
\item 反对称矩阵的性质：$[\omega_b]_{\times} \bm{v}_b = \omega_b \times \bm{v}_b$
\item 坐标变换：从惯性系到体坐标系需要左乘$R^T(t)$
\item $\omega_b \times \bm{v}_b$项是旋转坐标系效应的直接结果
\end{itemize}

\subsection{应用场景}

\begin{itemize}
\item 刚体动力学（欧拉方程）
\item 机器人动力学（递归牛顿-欧拉算法）
\item 航天器姿态控制
\item 旋转机械系统分析
\end{itemize}

\end{document}

